This appendix proves Proposition~\ref{prop:atom-morse}, the transversality
input used to transfer limiting active sets in
Theorem~\ref{thm:atom-count-limit}. The argument shows that every active KKT
contact of the limiting Poisson likelihood is nondegenerate and that no inactive contact
occurs. We condition on finitely many negative Poisson points and apply
parametric transversality; the elementary zero bound below supplies the
required Jacobian rank.

\begin{lemma}
\label{lem:atom-exp-polynomial}
Let \(\alpha_1,\ldots,\alpha_m\) be distinct real numbers, and let \(p_j\)
be real polynomials of degrees at most \(d_j\). Unless
\[
    A(x)=\sum_{j=1}^mp_j(x)e^{\alpha_jx}
\]
vanishes identically, it has at most
\(\sum_j(d_j+1)-1\) real zeros, counted with multiplicity.
\end{lemma}

\begin{proof}
The functions
\[
    e^{\alpha_jx},xe^{\alpha_jx},\ldots,x^{d_j}e^{\alpha_jx},
    \qquad 1\le j\le m,
\]
form an extended complete Chebyshev system on \(\mathbb R\). After their
common exponential factors are removed, the successive Wronskians are
nonzero confluent Vandermonde determinants. The repeated-Rolle argument for
such systems gives the stated zero bound; see
\cite[Chapter~I]{karlin1966tchebycheff}.
\end{proof}

We now prove the contact statement used in the main text. The continuously
distributed Poisson locations supply more perturbation parameters than any
bad-contact stratum has equations.
\begin{proof}[Proof of Proposition~\ref{prop:atom-morse}]
We show that each bad-contact event is null. Suppose first that the optimizer
has \(k\) distinct interior atoms with positive masses. Its parameters range
over the open stratum
\[
    s_j>0,
    \qquad b<\theta_1<\cdots<\theta_k<1.
\]
Choose \(q=6(k+1)^2\) negative Poisson points and isolate them in disjoint
intervals with rational endpoints, each containing exactly one point.
Conditional on the process outside these intervals and on the count events,
their locations \(x_1,\ldots,x_q\) have a joint density positive on the
product of the intervals. The rest of the process contributes a fixed
real-analytic background. The selected points contribute
\[
    \sum_{\ell=1}^q\log\{1+F_\nu(x_\ell)\}
\]
to the objective: inserting a point at \(x\) adds \(F_\nu(x)\) to the
linear field and \(h\{F_\nu(x)\}\) to the nonlinear term.

We work first on rational compact substrata on which masses are bounded away
from zero, atom locations are separated, and all free locations remain away
from their interval endpoints. These substrata countably cover every finite
atomic measure with distinct positive atoms.

The active parameters obey the \(2k\) analytic stationarity equations
\[
    \Gamma_\nu(\theta_j)=0,
    \qquad
    \Gamma_\nu'(\theta_j)=0,
    \qquad 1\le j\le k.
\]
An inactive interior contact \(\theta_0\) adds the two equations
\(\Gamma_\nu(\theta_0)=\Gamma_\nu'(\theta_0)=0\) and the one variable
\(\theta_0\). A flat active contact adds
\(\Gamma_\nu''(\theta_j)=0\) and no variable. Let \(\Phi\) collect the
active equations and the appropriate bad-contact equations.

We claim that \(\Phi\), viewed as a function of the atomic parameters, the
possible location \(\theta_0\), and \(x_1,\ldots,x_q\), is a submersion.
Every row of \(\Phi\) is a first variation in one of the generalized
directions
\[
    \delta_{\theta_j},\quad
    \delta_{\theta_j}',\quad
    \delta_{\theta_0},\quad
    \delta_{\theta_0}',\quad
    \delta_{\theta_j}''.
\]
Take a linear combination of these rows and denote the corresponding
generalized signed measure by \(\eta\). The contribution of a selected point
at \(x\) to this row combination is
\[
    \frac{F_\eta(x)}{1+F_\nu(x)}.
\]
If the combination annihilates every Jacobian column corresponding to a
selected location, then
\begin{equation}
\label{eq:atom-transversality-zeros}
    \frac d{dx}\left\{\frac{F_\eta(x)}{1+F_\nu(x)}\right\}=0
    \qquad\text{at }x=x_1,\ldots,x_q.
\end{equation}
After multiplication by \(\{1+F_\nu(x)\}^2\), the numerator is the
exponential polynomial
\[
    A_\eta(x)
    =F_\eta'(x)\{1+F_\nu(x)\}-F_\eta(x)F_\nu'(x).
\]
It uses at most \((k+1)^2\) distinct exponents, with polynomial
coefficients of degree at most two. Lemma~\ref{lem:atom-exp-polynomial}
bounds the number of zeros of a nonzero \(A_\eta\) by less than
\(3(k+1)^2\). Equation \eqref{eq:atom-transversality-zeros} supplies twice
that many, so \(A_\eta\equiv0\).

It follows that \(F_\eta/(1+F_\nu)\) is constant. Both numerator and
denominator tend respectively to zero and one as \(x\to-\infty\), so the
constant is zero and \(F_\eta\equiv0\). Linear independence of the
generalized exponentials gives \(\eta=0\). The generalized point masses
listed above are linearly independent because the active locations are
distinct and an inactive location is not active. Hence the Jacobian columns
in \(x_1,\ldots,x_q\) already have full row rank.

The parametric transversality theorem
\cite[Chapter~2]{guillemin2010differential} now applies. For almost every
selected-point configuration, the map obtained by fixing those locations is
transverse to zero. For an inactive interior contact, the domain has
dimension \(2k+1\) and the target has dimension \(2k+2\); for a flat active
contact, the corresponding dimensions are \(2k\) and \(2k+1\). In either
case the zero set is empty.

The boundary strata are analogous. If \(b\) is inactive, the bad equation
\(\Gamma_\nu(b)=0\) adds one equation and no variable. If \(b\) is active,
its location is fixed, so the active stratum has \(2k-1\) parameters, and
the bad equality \(\Gamma_\nu'(b)=0\) adds one equation. The same
generalized-direction argument, now using \(\delta_b\) or \(\delta_b'\),
makes both systems submersions with targets larger than their domains.

Finally remove the conditioning. Every locally finite simple Poisson
configuration contains arbitrarily many negative points that can be isolated
in disjoint rational intervals. There are only countably many choices of
\(k\), compact parameter substrata, contact windows, and rational interval
tuples, and every conditional bad event above has probability zero.
Colliding locations and zero masses require no extra stratum because the
optimizer is represented by its distinct support points with positive
masses.

Since \(\Gamma_\nu\le0\), an active interior zero has
\(\Gamma_\nu''\le0\), and an active boundary zero has
\(\Gamma_\nu'(b)\le0\). The null events just excluded are exactly the
equality cases. This proves all three assertions.
\end{proof}
